
\SGARefTarget{sga3:viii:section:6}\section*{6. Essentially Free Morphisms and Representability of Certain Functors of the Form \(\prod_{Y/S}Z/Y\)\textsuperscript{(*)}}
\addcontentsline{toc}{section}{6. Essentially Free Morphisms and Representability of Certain Functors}

\begingroup
\renewcommand{\thefootnote}{(*)}
\footnotetext[0]{\SGARefLink{sga3:viii:section:6}{This section} is independent of the theory of diagonalizable
groups; its natural place would be in Expos\'e VI\textsubscript{B}.}
\endgroup

\noindent\SGARefTarget{sga3:viii:definition:6:1}\textbf{Definition 6.1.}
Let \(f:X\to S\) be a morphism of preschemes. We say that \(f\) is
\emph{essentially free}, or that \(X\) is \emph{essentially free over
\(S\)}, if there are an affine open covering \((S_i)\) of \(S\), for
each \(i\) an \(S_i\)-prescheme \(S_i'\) that is affine and faithfully
flat over \(S_i\), and a covering \((X'_{ij})_j\) of
\(X_i'=X\times_SS_i'\) by affine opens \(X'_{ij}\), such that, for every
\((i,j)\), the ring of \(X'_{ij}\) is a free module over the ring of
\(S_i'\).\footnote{Editor's note: ``Free'' may be replaced by
``projective''; cf. \SGARefLink{sga3:viii:remark:6:8}{6.8} below. This
notion should also be compared with that of an \(S\)-prescheme that is
flat and pure, introduced and developed in \SGARefLink{sga3:viii:bib:rg71}{[RG71]}; see in particular
\emph{loc. cit.}, Part I, 3.3.12, and Part II, 3.1.4.1.}

\medskip
\noindent\SGARefTarget{sga3:viii:proposition:6:2}\textbf{Proposition 6.2.}
\begin{enumerate}
  \item[(a)]\SGARefTarget{sga3:viii:proposition:6:2:item:a:01}
  \textit{If \(X\) is essentially free over \(S\), then it is flat over
  \(S\). The converse holds if \(S\) is Artinian.}

  \item[(b)]\SGARefTarget{sga3:viii:proposition:6:2:item:b:01}
  \textit{If \(S\) is the spectrum of a field, every \(S\)-prescheme is
  essentially free over \(S\).}

  \item[(c)]\SGARefTarget{sga3:viii:proposition:6:2:item:c:01}
  \textit{If \(X\) is essentially free over \(S\), and if \(S'\to S\)
  is a base-change morphism, then \(X'=X\times_SS'\) is essentially
  free over \(S'\). The converse holds if \(S'\to S\) is faithfully
  flat and quasi-compact.}
\end{enumerate}

The proof is immediate, using for the converse in \SGARefLink{sga3:viii:proposition:6:2:item:a:01}{(a)} the fact that a
flat module over an Artinian local ring is free.\footnote{Editor's note:
Indeed, let \((A,\mathfrak m)\) be an Artinian local ring, \(k\) its
residue field, \(M\) an arbitrary \(A\)-module, and
\((x_i)_{i\in I}\) elements of \(M\) whose images form a basis of
\(M/\mathfrak mM\) over \(k\). Let \(F\) be the free \(A\)-module with
basis \((e_i)_{i\in I}\), and let \(\phi:F\to M\) be the
\(A\)-morphism defined by \(\phi(e_i)=x_i\). Then
\(Q=\operatorname{Coker}\phi\) satisfies \(Q=\mathfrak mQ\), whence
\(Q=0\), since \(\mathfrak m\) is nilpotent. Suppose moreover that
\(M\) is flat over \(A\). Then \(K=\operatorname{Ker}\phi\) satisfies
\(K\otimes_Ak=0\), that is, \(K=\mathfrak mK\), whence \(K=0\).}

\medskip
\noindent\SGARefTarget{sga3:viii:proposition:6:3}\textbf{Proposition 6.3.}
\textit{Let \(H\) be a diagonalizable \(S\)-group prescheme
(more generally, one that becomes diagonalizable after a suitable
faithfully flat quasi-compact extension of every affine open of \(S\);
that is, \(H\) is ``of multiplicative type,'' cf. IX, \SGARefLink{sga3:ix:definition:1:1}{1.1}). Then \(H\)
is essentially free over \(S\).}

Indeed, if \(H\) is diagonalizable, it is affine over \(S\) and is
defined by an algebra that is a free \(\mathcal O_S\)-module.

The introduction of \SGARefLink{sga3:viii:definition:6:1}{Definition 6.1}
is justified by \SGARefLink{sga3:viii:theorem:6:4}{the following theorem}.

\medskip
\noindent\SGARefTarget{sga3:viii:theorem:6:4}\textbf{Theorem 6.4.}
\textit{Let \(S\) be a prescheme, \(Z\) an \(S\)-prescheme essentially
free over \(S\), and \(Y\) a closed subprescheme of \(Z\). Consider the
functor}\footnote{Editor's note: Cf. II.1, where \SGARefLink{sga3:viii:theorem:6:4:display:functor:01}{this functor} is denoted
\(\prod_{Z/S}Y\).}
\SGARefTarget{sga3:viii:theorem:6:4:display:functor:01}
\[
  \begin{aligned}
  F=\prod_{Z/S}Y/Z:\;(\mathbf{Sch}/S)^\circ
    &\longrightarrow\mathbf{Set},\\
  F(S')=\Gamma(Y_{S'}/Z_{S'})
    &=
    \begin{cases}
      \varnothing,& Z_{S'}\ne Y_{S'},\\
      \{\operatorname{id}_{Z_{S'}}\},& Z_{S'}=Y_{S'}.
    \end{cases}
  \end{aligned}
\]
\textit{\SGARefLink{sga3:viii:theorem:6:4:display:functor:01}{This functor} is representable by a closed subprescheme of
\(S\).}\footnote{Editor's note: We have corrected \(Z\) to \(S\).}

\smallskip
\noindent\footnote{Editor's note: We have expanded the original in what
follows; see also \SGARefLink{sga3:viii:corollary:1:7:3}{1.7.3}.}
First note that \(F\) is a sheaf for the \((\mathrm{fpqc})\) topology.
Since \(F(S')\) is either \(\varnothing\) or \(\{\mathrm{pt}\}\) for
every \(S'\), it is enough to check the following assertion: if
\((S_i)\) is an open covering of \(S\) (respectively, if
\(S'\to S\) is faithfully flat and quasi-compact), and every
\(Y_{S_i}\to Z_{S_i}\) is an isomorphism (respectively,
\(Y_{S'}\to Z_{S'}\) is an isomorphism), then \(Y\to Z\) is an
isomorphism. This is clear in the first case and follows from SGA 1,
VIII, 5.4, or EGA IV\textsubscript{2}, 2.7.1, in the second.

Moreover, by SGA 1, VIII, 2.1 and 5.5, faithfully flat quasi-compact
morphisms are morphisms of effective descent for the fibered category
of arrows that are closed immersions. Thus, with the notation of
\SGARefLink{sga3:viii:definition:6:1}{6.1}, we may reduce to the case
\(S=S_i'\).

Let \((Z_j)\) be a covering of \(Z\) by affine opens such that
\(\mathcal O(Z_j)\) is a free module over \(A=\mathcal O(S)\), put
\(Y_j=Y\cap Z_j\), and let
\[
  F_j:(\mathbf{Sch}/S)^\circ\longrightarrow\mathbf{Set}
\]
be the functor defined from \((Z_j,Y_j)\) as \(F\) is defined from
\((Z,Y)\). It is a subfunctor of the terminal functor, and clearly
\[
  F=\bigcap_jF_j.
\]

It therefore suffices to prove that every \(F_j\) is representable by a
closed subscheme \(T_j\) of \(S\), for then \(F\) is represented by the
closed subscheme \(T=\bigcap_jT_j\). We may consequently suppose that
\(Z\) is affine as well, say \(Z=\operatorname{Spec}(B)\), where \(B\)
is a free \(A\)-module. Let \(J\) be a subset of \(B\) defining the
closed subprescheme \(Y\) of \(Z\), and let \(K\) be the ideal of \(A\)
generated by the \(u_i(J)\subseteq A\), where the \(u_i:B\to A\) are
the coordinate forms for the chosen basis. It is immediate that
\[
  T=\mathcal V(K)=\operatorname{Spec}(A/K)
\]
has the required property. This completes the proof.

\medskip
\noindent\SGARefTarget{sga3:viii:examples:6:5}\textbf{Examples 6.5.}
We give some important examples of functors that reduce to functors
\(\prod_{Z/S}Y/Z\) of the type considered in
\SGARefLink{sga3:viii:theorem:6:4}{6.4}, and for which criteria of
representability will be useful below. Here \(S\) denotes a prescheme,
and \(X,Y,Z,\ldots\) denote preschemes over \(S\).

\smallskip
\noindent\SGARefTarget{sga3:viii:examples:6:5:item:a:01}\textit{(a)}
Suppose given an \(S\)-morphism
\[
  q:X\longrightarrow\underline{\operatorname{Hom}}_S(Y,Z),
  \tag{x}\label{sga3:viii:examples:6:5:eq:x:01}
\]
(``\(X\) acts on \(Y\) with values in \(Z\)''), that is, a morphism
\[
  r:X\times_SY\longrightarrow Z.
  \tag{xx}\label{sga3:viii:examples:6:5:eq:xx:01}
\]
Let \(Z'\) be a subprescheme of \(Z\). It gives a monomorphism
\[
  \underline{\operatorname{Hom}}_S(Y,Z')
  \longrightarrow
  \underline{\operatorname{Hom}}_S(Y,Z),
\]
making the first functor a subfunctor of the second. Let \(X'\) be the
inverse image of this subfunctor under
\SGARefLink{sga3:viii:examples:6:5:eq:x:01}{(x)}. Thus \(X'\) is the
subfunctor of \(X\) for which \(X'(T)\) is the set of \(x\in X(T)\)
such that
\[
  q(x):Y_T\longrightarrow Z_T
\]
factors through \(Z_T'\). This functor has the following description.
Put \(P=X\times_SY\), and let \(P'\) be the inverse image of \(Z'\)
under \(r:P\to Z\). Then there is an evident isomorphism
\[
  X'\simeq\prod_{P/X}P'/P.
  \tag{xxx}\label{sga3:viii:examples:6:5:eq:xxx:01}
\]
It follows that \textit{if \(Y\) is essentially free over \(S\) and
\(Z'\) is closed in \(Z\), then the subfunctor \(X'\) of \(X\) is
representable by a closed subprescheme of \(X\).}

\smallskip
\noindent\SGARefTarget{sga3:viii:examples:6:5:item:b:01}\textit{(b)}
Suppose given two actions of \(X\) on \(Y\) with values in \(Z\), that
is, two morphisms
\[
  q_1,q_2:X\rightrightarrows
  \underline{\operatorname{Hom}}_S(Y,Z),
\]
and put \(X'=\operatorname{Ker}(q_1,q_2)\). This is the subfunctor of
\(X\) for which \(X'(T)\) is the set of \(x\in X(T)\) such that the two
morphisms
\[
  q_1(x),q_2(x):Y_T\rightrightarrows Z_T
\]
are equal. Giving \(q_1,q_2\) is equivalent to giving a morphism
\[
  q:X\longrightarrow
  \underline{\operatorname{Hom}}_S(Y,Z\times_SZ),
\]
or, equivalently, a morphism
\[
  r:X\times_SY\longrightarrow Z\times_SZ.
\]
Put \(U=Z\times_SZ\), and let \(U'\) be the diagonal subprescheme of
\(Z\times_SZ\). Then \(X'\) is the inverse image under \(q\) of the
subfunctor
\[
  \underline{\operatorname{Hom}}_S(Y,U')
  \longrightarrow
  \underline{\operatorname{Hom}}_S(Y,U).
\]
It can therefore be written in the form
\SGARefLink{sga3:viii:examples:6:5:eq:xxx:01}{(xxx)}, with
\(P=X\times_SY\) and \(P'\) the inverse image of the diagonal under
\(r\), that is, the kernel of
\[
  X\times_SY
  \mathop{\rightrightarrows}^{r_1}_{r_2}
  Z.
\]
Thus we are in the situation of \SGARefLink{sga3:viii:examples:6:5:item:a:01}{(a)}.

Consequently, \textit{if \(Y\) is essentially free over \(S\) and \(Z\)
is separated over \(S\), then the subfunctor \(X'\) of \(X\) is
representable by a closed subprescheme of \(X\).}

\smallskip
\noindent\SGARefTarget{sga3:viii:examples:6:5:item:c:01}\textit{(c)}
Suppose given a morphism
\SGARefTarget{sga3:viii:examples:6:5:item:c:display:q:01}\[
  q:X\longrightarrow\underline{\operatorname{Hom}}_S(Y,Y);
\]
that is, ``\(X\) acts on \(Y\).'' Let \(X'\) be the ``kernel'' of \SGARefLink{sga3:viii:examples:6:5:item:c:display:q:01}{this
morphism}, namely the subfunctor of \(X\) for which \(X'(T)\) is the set
of \(x\in X(T)\) such that \(q(x):Y_T\to Y_T\) is the identity. This
functor falls under \SGARefLink{sga3:viii:examples:6:5:item:b:01}{(b)}, as one sees by introducing a second morphism
\[
  q':X\longrightarrow\underline{\operatorname{Hom}}_S(Y,Y)
\]
that ``makes \(X\) act trivially on \(Y\).'' Hence, \textit{if \(Y\) is
essentially free and separated over \(S\), then the kernel subfunctor of
\(q\) is representable by a closed subprescheme of \(X\).}

\smallskip
\noindent\SGARefTarget{sga3:viii:examples:6:5:item:d:01}\textit{(d)}
Under the hypotheses of \SGARefLink{sga3:viii:examples:6:5:item:c:01}{(c)}, consider the subfunctor \(Y'\) of \(Y\) of
``invariants under \(X\).'' Thus \(Y'(T)\) is the set of \(y\in Y(T)\)
such that the corresponding morphism \(\bar q(y):X_T\to Y_T\) is the
``\(T\)-constant morphism with value \(y\).'' Introducing \(q'\) as in
\SGARefLink{sga3:viii:examples:6:5:item:c:01}{(c)}, and writing \(\bar q\) and \(\bar q'\) for the corresponding morphisms,
\[
  \bar q,\bar q':Y\rightrightarrows
  \underline{\operatorname{Hom}}_S(X,Y),
\]
one sees that \(Y'=\operatorname{Ker}(\bar q,\bar q')\). It therefore again
falls under \SGARefLink{sga3:viii:examples:6:5:item:b:01}{(b)}, with the roles of \(X\) and \(Y\) reversed and \(Z=Y\).
Consequently, \textit{if \(X\) is essentially free over \(S\) and \(Y\)
is separated over \(S\), then the subfunctor \(Y'\) of \(Y\) of
invariants under \(X\) is representable by a closed subprescheme of
\(Y\).}

\smallskip
\noindent\SGARefTarget{sga3:viii:examples:6:5:item:e:01}\textit{(e)}
Constructions of the type described in \SGARefLink{sga3:viii:examples:6:5}{the preceding examples} occur
especially often in group theory. Thus, when a group prescheme \(G\)
over \(S\) acts on the \(S\)-prescheme \(X\),
\[
  q:G\longrightarrow\underline{\operatorname{Aut}}_S(X),
\]
the kernel of \(q\), the ``subgroup of \(G\) acting trivially,'' is a
closed subprescheme of \(G\) provided that \(X\) is essentially free
and separated over \(S\), by \SGARefLink{sga3:viii:examples:6:5:item:c:01}{(c)}. The subobject \(X^G\) of invariants
is a closed subprescheme of \(X\), provided that \(G\) is essentially
free over \(S\) and \(X\) is separated over
\(S\),\footnote{Editor's note: We have corrected \(S_X\) to \(S\).}
by \SGARefLink{sga3:viii:examples:6:5:item:d:01}{(d)}.

Let \(Y,Z\) be subpreschemes of \(X\), and consider the subfunctor
\(\underline{\operatorname{Transp}}_G(Y,Z)\) of \(G\), the
``transporter of \(Y\) into \(Z\).'' Its points with values in an
\(S\)-prescheme \(T\) are the \(g\in G(T)\) such that the corresponding
automorphism of \(X_T\) satisfies \(g(Y_T)\subseteq Z_T\), that is,
induces a morphism \(Y_T\to X_T\) factoring through \(Z_T\). Hence,
\textit{if \(Y\) is essentially free over \(S\), and \(Z\) is closed in
\(X\), then \(\underline{\operatorname{Transp}}_G(Y,Z)\) is a closed
subprescheme of \(G\)}, by \SGARefLink{sga3:viii:examples:6:5:item:a:01}{(a)}.

One may also consider the \emph{strict transporter} of \(Y\) into
\(Z\),\footnote{Editor's note: It is denoted
\(\underline{\operatorname{Transp}}^{\mathrm{str}}_G(Y,Z)\).}
whose \(T\)-valued points are the \(g\in G(T)\) such that
\(g(Y_T)=Z_T\). It is precisely
\[
  \underline{\operatorname{Transp}}_G(Y,Z)
  \cap
  \sigma\bigl(\underline{\operatorname{Transp}}_G(Z,Y)\bigr),
\]
where \(\sigma\) is the symmetry of \(G\). Consequently, if \(Y\) and
\(Z\) are essentially free over \(S\) and closed in \(X\), then the
strict transporter of \(Y\) into \(Z\) is a closed subprescheme of
\(G\).

An important case is \(X=G\), with \(G\) acting on itself by inner
automorphisms. If \(H\) is a subprescheme of \(G\), the strict
transporter of \(H\) into itself is also called the \emph{normalizer of
\(H\) in \(G\)}, and is denoted
\(\underline{\operatorname{Norm}}_G H\). Hence, \textit{if \(H\) is a
closed subgroup prescheme of \(G\), essentially free over \(S\), then
\(\underline{\operatorname{Norm}}_G H\) is representable by a closed
subgroup prescheme of \(G\).}

Finally, let \(Z\) be a subprescheme of \(G\). Its \emph{centralizer}
\(\underline{\operatorname{Centr}}_G(Z)\) in \(G\) is the group
subfunctor of \(G\) defined by the procedure in \SGARefLink{sga3:viii:examples:6:5:item:d:01}{(d)}, regarding ``\(Z\)
acts on \(G\)'' through the action induced by that of \(G\). Thus,
\textit{if \(Z\) is essentially free over \(S\) and \(G\) is separated
over \(S\), then \(\underline{\operatorname{Centr}}_G(Z)\) is a closed
subgroup prescheme of \(G\).} In particular, \textit{if \(G\) is
essentially free and separated over \(S\), then the center \(C\) of
\(G\), which is none other than
\(\underline{\operatorname{Centr}}_G(G)\), is a closed subgroup
prescheme of \(G\).}

When \(S\) is the spectrum of a field,
\SGARefLink{sga3:viii:proposition:6:2:item:b:01}{6.2(b)}
\emph{[the source prints 6.3(b)]} shows that in examples
\SGARefLink{sga3:viii:examples:6:5:item:a:01}{(a)}--\SGARefLink{sga3:viii:examples:6:5:item:e:01}{(e)} above the ``essentially free'' hypotheses are automatically
satisfied; only separation hypotheses remain. Recalling that a
group prescheme over a field is necessarily separated, one obtains,
for example:

\medskip
\noindent\SGARefTarget{sga3:viii:corollary:6:7}\textbf{Corollary 6.7.}
\footnote{Editor's note: We have retained the numbering of the original;
there is no 6.6.}
\textit{Let \(G\) be a group prescheme over a field \(k\). Then:}
\begin{list}{--}{\setlength{\leftmargin}{1.8em}\setlength{\itemsep}{0.35em}}
  \item\SGARefTarget{sga3:viii:corollary:6:7:item:1:01}
  \textit{For every subprescheme \(Z\) of \(G\), the centralizer of
  \(Z\) in \(G\) is a closed subgroup prescheme of \(G\). In particular,
  this holds for the center
  \(\underline{\operatorname{Centr}}_G(G)\) of \(G\).}

  \item\SGARefTarget{sga3:viii:corollary:6:7:item:2:01}
  \textit{More generally, if \(u,v:X\to G\) are morphisms of
  preschemes, then \(\underline{\operatorname{Transp}}_G(u,v)\) is
  representable by a closed subprescheme of \(G\).}

  \item\SGARefTarget{sga3:viii:corollary:6:7:item:3:01}
  \textit{For two subpreschemes \(Y,Z\) of \(G\), with \(Z\) closed,
  \(\underline{\operatorname{Transp}}_G(Y,Z)\) is a closed subprescheme
  of \(G\). If \(Y\) is also closed, \SGARefLink{sga3:viii:corollary:6:7:item:3:01}{the same conclusion} holds for
  \(\underline{\operatorname{Transp}}^{\mathrm{str}}_G(Y,Z)\).}

  \item\SGARefTarget{sga3:viii:corollary:6:7:item:4:01}
  \textit{For every subgroup prescheme}\footnote{Editor's note: Indeed,
  over a field \(k\), every subgroup prescheme of \(G\) is closed; cf.
  Expos\'e VI\textsubscript{A}, 0.6.1.}
  \textit{\(H\) of \(G\), the normalizer
\(\underline{\operatorname{Norm}}_G(H)\) is a closed subgroup
prescheme of \(G\).}
\end{list}

\medskip
\noindent\SGARefTarget{sga3:viii:remark:6:8}\textbf{Remark 6.8.}
\footnote{Editor's note: We have expanded the original in what follows;
in particular, we have added \SGARefLink{sga3:viii:lemma:6:8:1}{Lemma 6.8.1}.}
Let \(A\) be a commutative ring, let \(M\) be an \(A\)-module, and put
\[
  M^\vee=\operatorname{Hom}_A(M,A).
\]
Endow the \(A\)-module \(\operatorname{End}_A(M)\) with the topology of
discrete pointwise convergence. Thus a basis of neighborhoods of zero
is formed by the following \(A\)-submodules, where \(n\in\mathbb N\)
and \(m_1,\ldots,m_n\in M\):
\SGARefTarget{sga3:viii:remark:6:8:display:neighborhood:01}
\[
  K(m_1,\ldots,m_n)
  =
  \{u\in\operatorname{End}_A(M)\mid
  u(m_i)=0\ \text{for \(i=1,\ldots,n\)}\}.
\]
We say that \(M\) is a \emph{quasi-free \(A\)-module} if the image of
the canonical morphism
\[
  \Theta:M\otimes_AM^\vee\longrightarrow\operatorname{End}_A(M)
\]
contains \(\operatorname{id}_M\) in its closure; that is, if \SGARefLink{sga3:viii:remark:6:8:eq:star:01}{the
following condition} holds:

\[
  \left\{
  \begin{aligned}
    &\text{for all \(m_1,\ldots,m_n\in M\), there exist
      \(x_1,\ldots,x_r\in M\) and \(f_1,\ldots,f_r\in M^\vee\)}\\
    &\text{such that }\quad
    m_i
    =
    \Theta\!\left(\sum_{s=1}^r x_s\otimes f_s\right)(m_i)
    =
    \sum_{s=1}^r f_s(m_i)x_s
    \quad\text{for \(i=1,\ldots,n\).}
  \end{aligned}
  \right.
  \tag{*}\label{sga3:viii:remark:6:8:eq:star:01}
\]
\noindent(In this case \(\operatorname{Im}\Theta\) is dense in
\(\operatorname{End}_A(M)\), since for every
\(u\in\operatorname{End}_A(M)\) one has
\(
  u(m_i)
  =
  \sum_{s=1}^r f_s(m_i)u(x_s)
  =
  \Theta\!\left(\sum_{s=1}^r u(x_s)\otimes f_s\right)(m_i).
\))

First note that \SGARefLink{sga3:viii:remark:6:8:eq:star:01}{this property} is stable under base change. Indeed, let
\(\phi:A\to A'\) be a ring homomorphism, put \(M'=M\otimes_AA'\), and
let \(m_1',\ldots,m_n'\in M'\). Write
\[
  m_i'=\sum_jm_{ij}\otimes b_{ij},
  \qquad m_{ij}\in M,\quad b_{ij}\in A'.
\]
By hypothesis there exist \(x_1,\ldots,x_r\in M\) and
\(f_1,\ldots,f_r\in\operatorname{Hom}_A(M,A)\) such that
\[
  m_{ij}=\sum_sx_sf_s(m_{ij})
  \qquad\text{for all \(i,j\).}
\]
Write \(\phi\circ f_s\) for the image of \(f_s\) in
\[
  \operatorname{Hom}_A(M,A')
  =
  \operatorname{Hom}_{A'}(M',A').
\]
Then, for \(i=1,\ldots,n\),
\SGARefTarget{sga3:viii:remark:6:8:display:base-change:01}
\[
  \begin{aligned}
  \left(\sum_sx_s\otimes\phi\circ f_s\right)(m_i')
    &=
    \sum_{s,j}x_s\otimes\phi(f_s(m_{ij}))\,b_{ij}\\
    &=
    \sum_j\left(\sum_sx_sf_s(m_{ij})\right)\otimes b_{ij}\\
    &=m_i'.
  \end{aligned}
\]
This proves that \(M'\) is quasi-free over \(A'\).

Also, every projective \(A\)-module \(P\) is quasi-free. Indeed, there
are \(A\)-morphisms
\[
  A^{(I)}\xrightarrow{\ \pi\ }P
  \xrightarrow{\ \tau\ }A^{(I)}
  \qquad\text{such that}\qquad
  \pi\circ\tau=\operatorname{id}_P.
\]
Let \((e_i)\) be the canonical basis of \(A^{(I)}\), and put
\(f_i=e_i^*\circ\tau\), a linear form on \(P\). Given
\(m_1,\ldots,m_n\in P\), there is a finite subset \(J\) of \(I\) such
that
\[
  m_k=\sum_{i\in J}f_i(m_k)\pi(e_i)
  \qquad\text{for \(k=1,\ldots,n\).}
\]

Theorem \SGARefLink{sga3:viii:theorem:6:4}{6.4} remains valid if, in the
statement of \SGARefLink{sga3:viii:definition:6:1}{Definition 6.1},
``free'' is replaced by ``projective'' or, more generally, by
``quasi-free.'' Proceeding as in the proof of
\SGARefLink{sga3:viii:theorem:6:4}{6.4}, one reduces to \SGARefLink{sga3:viii:lemma:6:8:1}{the following
lemma}.

\medskip
\noindent\SGARefTarget{sga3:viii:lemma:6:8:1}\textbf{Lemma 6.8.1.}
\textit{Let \(M\) be a quasi-free \(A\)-module, \(N\) a submodule, and
\(F\) the covariant functor}
\[
  F:(\text{\(A\)-algebras})\longrightarrow\mathbf{Set}
\]
\textit{such that \(F(B)=\{\mathrm{pt}\}\) if
\(M_B=(M/N)_B\), and \(F(B)=\varnothing\) otherwise. Then there is an
ideal \(K\) of \(A\) such that \(F(B)=\{\mathrm{pt}\}\) if and only if
the morphism \(A\to B\) factors through \(A/K\); that is, there is an
isomorphism, functorial in \(B\),}
\SGARefTarget{sga3:viii:lemma:6:8:1:display:representing:01}
\[
  F(B)\simeq\operatorname{Hom}_{A\text{-alg.}}(A/K,B).
\]

\smallskip
\noindent\textit{Proof.}
Let \((n_j)\) be a system of generators of \(N\), and let \(F_j\) be the
subfunctor of the terminal functor \(e\), where
\(e(B)=\{\mathrm{pt}\}\) for every \(B\), corresponding to the
submodule \(An_j\). We have \(F(B)=\{\mathrm{pt}\}\) if and only if the
image of every \(n_j\) in \(M_B\) is zero. Hence \(F\) is the
intersection of the functors \(F_j\), reducing us to the case in which
\(N\) is generated by one element \(n\).

Let \(B\) be an \(A\)-algebra, with structure map \(\phi:A\to B\).
If the image \(n\otimes1\) of
\(n\) in \(M_B\) is zero, then for every
\(f\in M^\vee=\operatorname{Hom}_A(M,A)\),
\[
  0=f(n)\otimes1=\phi(f(n)).
\]
Conversely, because \(M\) is quasi-free, there exist
\(x_1,\ldots,x_r\in M\) and \(f_1,\ldots,f_r\in M^\vee\) such that
\[
  n=\sum_sf_s(n)x_s,
\]
whence
\[
  n\otimes1=\sum_sx_s\otimes\phi(f_s(n)).
\]
It follows that \(n\otimes1=0\) if and only if \(\phi\) factors through
\(A/K\), where \(K\) is the ideal generated by the \(f_s(n)\), for
\(s=1,\ldots,r\). This proves the lemma.
